% ═══════════════════════════════════════════════════════════════
% LH-01d: ¿Qué tal? + Stundenleistung (LEHRERHINWEISE)
% Spanisch Klasse 7 - Sequenz 1: Empezamos (ABSCHLUSS)
% ═══════════════════════════════════════════════════════════════

\documentclass[11pt, a4paper]{article}

\newif\ifswmodus
\swmodusfalse

\usepackage[utf8]{inputenc}
\usepackage[T1]{fontenc}
\usepackage[ngerman]{babel}
\usepackage{helvet}
\renewcommand{\familydefault}{\sfdefault}

\usepackage[margin=2cm]{geometry}
\usepackage{enumitem}
\usepackage{tabularx}
\usepackage{booktabs}
\usepackage{longtable}

\usepackage{xcolor}
\usepackage{amssymb}
\usepackage{tcolorbox}
\tcbuselibrary{skins}

\usepackage{fancyhdr}
\usepackage{lastpage}
\usepackage{needspace}

% FARBEN
\definecolor{spanisch-color}{HTML}{c41e3a}
\definecolor{grau-color}{HTML}{888888}
\definecolor{hellgrau-color}{HTML}{f5f5f5}
\definecolor{basis-color}{HTML}{2a7a4b}
\definecolor{standard-color}{HTML}{2c5aa0}
\definecolor{erweiterung-color}{HTML}{7b1fa2}
\definecolor{loesung-color}{HTML}{c62828}

\definecolor{sw-dark}{HTML}{000000}
\definecolor{sw-gray}{HTML}{666666}
\definecolor{sw-white}{HTML}{ffffff}

\ifswmodus
    \colorlet{spanisch}{sw-dark}
    \colorlet{grau}{sw-gray}
    \colorlet{hellgrau}{sw-white}
    \colorlet{basis}{sw-dark}
    \colorlet{standard}{sw-dark}
    \colorlet{erweiterung}{sw-dark}
    \colorlet{loesung}{sw-dark}
\else
    \colorlet{spanisch}{spanisch-color}
    \colorlet{grau}{grau-color}
    \colorlet{hellgrau}{hellgrau-color}
    \colorlet{basis}{basis-color}
    \colorlet{standard}{standard-color}
    \colorlet{erweiterung}{erweiterung-color}
    \colorlet{loesung}{loesung-color}
\fi

\newcommand{\fachfarbe}{spanisch}

\newcommand{\thema}{¿Qué tal? + Stundenleistung}
\newcommand{\sequenz}{01}
\newcommand{\block}{d}
\newcommand{\fach}{Spanisch}
\newcommand{\klasse}{7}
\newcommand{\dauer}{90 Min. (Doppelstunde)}
\newcommand{\lerngruppengroesse}{24}
\newcommand{\datum}{\today}

\pagestyle{fancy}
\fancyhf{}
\renewcommand{\headrulewidth}{0.4pt}
\renewcommand{\footrulewidth}{0.4pt}

\lhead{\fach{} -- Klasse \klasse}
\chead{\textbf{LH-\sequenz\block{} (Lehrerhinweise)}}
\rhead{\datum}

\lfoot{}
\cfoot{}
\rfoot{Seite \thepage\ von \pageref{LastPage}}

\newcommand{\B}{\textcolor{basis}{\textbf{[B]}}}
\newcommand{\Std}{\textcolor{standard}{\textbf{[S]}}}
\newcommand{\E}{\textcolor{erweiterung}{\textbf{[E]}}}
\newcommand{\loesung}[1]{\textcolor{loesung}{\textbf{#1}}}
\newcommand{\es}[1]{\textcolor{\fachfarbe}{\textbf{#1}}}

\newtcolorbox{kurzplan}{
    colback=hellgrau,
    colframe=\fachfarbe,
    colbacktitle=white,
    coltitle=\fachfarbe,
    boxrule=1pt,
    arc=0pt,
    fonttitle=\bfseries,
    attach boxed title to top left={yshift=-2mm, xshift=5mm},
    boxed title style={boxrule=0pt, colback=white},
    title=Kurzplan
}

\newtcolorbox{differenzierung}{
    colback=white,
    colframe=grau,
    colbacktitle=white,
    coltitle=grau,
    boxrule=0.5pt,
    arc=0pt,
    fonttitle=\bfseries,
    attach boxed title to top left={yshift=-2mm, xshift=5mm},
    boxed title style={boxrule=0pt, colback=white},
    title=Differenzierung
}

\newtcolorbox{fehlerbox}{
    colback=white,
    colframe=loesung,
    colbacktitle=white,
    coltitle=loesung,
    boxrule=0.5pt,
    arc=0pt,
    fonttitle=\bfseries,
    attach boxed title to top left={yshift=-2mm, xshift=5mm},
    boxed title style={boxrule=0pt, colback=white},
    title=Häufige Fehler
}

\newtcolorbox{materialbox}{
    colback=white,
    colframe=\fachfarbe,
    colbacktitle=white,
    coltitle=\fachfarbe,
    boxrule=0.5pt,
    arc=0pt,
    fonttitle=\bfseries,
    attach boxed title to top left={yshift=-2mm, xshift=5mm},
    boxed title style={boxrule=0pt, colback=white},
    title=Material-Checkliste
}

\begin{document}

\begin{center}
    {\Large\bfseries\textcolor{\fachfarbe}{Lehrerhinweise: \thema}}\\[0.3em]
    {\normalsize LH-\sequenz\block{} | \dauer{} | \textbf{ABSCHLUSS Sequenz 1}}
\end{center}

\vspace{10pt}

\section*{1. Stundenverlauf (Kurzplan)}

\textbf{1. Stunde (45 Min.) -- Schwerpunkt: ¿Qué tal? + Übung}

\begin{tabularx}{\textwidth}{|c|l|X|l|}
\hline
\textbf{Zeit} & \textbf{Phase} & \textbf{Inhalt} & \textbf{Material} \\
\hline
0--5 & Einstieg & Begrüßung mit ``¡Hola! ¿Qué tal?'' & -- \\
\hline
5--15 & Input & ¿Qué tal? + Antworten einführen & Tafel \\
\hline
15--25 & Übung 1 & AB-01d Aufg. 1 (Zuordnung) & AB \\
\hline
25--35 & Übung 2 & Dialog-Kette: Reihum fragen/antworten & -- \\
\hline
35--45 & Sicherung & Merkkasten ausfüllen & AB \\
\hline
\end{tabularx}

\vspace{1em}

\textbf{2. Stunde (45 Min.) -- Schwerpunkt: Stundenleistung}

\begin{tabularx}{\textwidth}{|c|l|X|l|}
\hline
\textbf{Zeit} & \textbf{Phase} & \textbf{Inhalt} & \textbf{Material} \\
\hline
0--5 & Aktivierung & Schnelle Wiederholung: ser, Artikel & -- \\
\hline
5--10 & Vorbereitung & SL erklären, Fragen klären & SL \\
\hline
10--35 & \textbf{TEST} & Stundenleistung (25 Min.) & SL \\
\hline
35--40 & Abgabe & Einsammeln & -- \\
\hline
40--45 & Abschluss & Ausblick Sequenz 2, Feedback & -- \\
\hline
\end{tabularx}

\section*{2. Differenzierung}

\begin{differenzierung}
\textbf{Legende:} \B{} Basis (BR) | \Std{} Standard (MR) | \E{} Erweiterung (GY)

\textbf{WICHTIG:} Keine Niveaumarkierungen auf Schülermaterial!
\end{differenzierung}

\vspace{1em}

\begin{tabularx}{\textwidth}{|l|X|X|X|}
\hline
& \B{} \textbf{Basis} & \Std{} \textbf{Standard} & \E{} \textbf{Erweiterung} \\
\hline
\textbf{AB-01d} & Aufg. 1-3 Pflicht & Alle Aufgaben & + freier Dialog \\
\hline
\textbf{SL-01d} & Aufg. 1-3 (12P = 4) & Alle (20P) & Bonusaufgabe \\
\hline
\textbf{Zeit SL} & +5 Min. Zugabe & 25 Min. & 25 Min. \\
\hline
\end{tabularx}

\vspace{1em}

\textbf{Konkrete Hinweise:}
\begin{itemize}
    \item \B{} SuS mit Prüfungsangst: Vorher kurz ermutigen, Basisaufgaben zuerst
    \item \B{} 12/20 Punkte = Note 4 (ausreichend)
    \item \E{} Schnelle: Dürfen Bonusaufgabe bearbeiten (nicht in Punktzahl)
\end{itemize}

\section*{3. Häufige Fehler}

\begin{fehlerbox}
\begin{tabularx}{\textwidth}{|l|X|X|}
\hline
\textbf{Fehler} & \textbf{Beispiel} & \textbf{Reaktion} \\
\hline
¿Qué tal? = Name & ``Me llamo bien'' & Unterschied klären: tal = Befinden \\
\hline
ser statt estar & *``Soy bien'' statt ``Estoy bien'' & Für Kl. 7: ``Bien'' ohne Verb akzeptieren \\
\hline
Artikel vergessen & *``libro'' statt ``el libro'' & Genus-Regel wiederholen \\
\hline
ser-Endungen & *``yo soes'' & Konjugationstabelle zeigen \\
\hline
\end{tabularx}
\end{fehlerbox}

\section*{4. Material-Checkliste}

\begin{materialbox}
\begin{itemize}[label=$\square$]
    \item AB-\sequenz\block.pdf (\lerngruppengroesse{} Kopien)
    \item ML-\sequenz\block.pdf (1× für Lehrkraft)
    \item \textbf{SL-\sequenz\block.pdf (\lerngruppengroesse{} Kopien) -- TEST!}
    \item Lehrwerk: Qué pasa 1, S. 22--27
    \item Tafelmarker / Kreide
    \item Stoppuhr/Timer für Stundenleistung
\end{itemize}
\end{materialbox}

\section*{5. Lösungen AB-\sequenz\block}

\textbf{Merkkasten:}
\begin{itemize}
    \item ¿Qué tal? = \loesung{Wie geht's?}
    \item ¿Cómo estás? = \loesung{Wie geht es dir?}
    \item Muy bien = \loesung{Sehr gut} | Bien = \loesung{Gut}
    \item Regular = \loesung{Es geht so} | Mal = \loesung{Schlecht}
    \item ¿Y tú? = \loesung{Und du?}
\end{itemize}

\textbf{Aufgabe 1 (Zuordnung):} (4P)
\begin{itemize}
    \item Muy bien $\rightarrow$ \loesung{++}
    \item Bien $\rightarrow$ \loesung{+}
    \item Regular $\rightarrow$ \loesung{0}
    \item Mal $\rightarrow$ \loesung{--}
\end{itemize}

\textbf{Aufgabe 2 (Dialog):} (6P)
\begin{itemize}
    \item A: ¡Hola! \loesung{¿Qué tal?}
    \item B: Bien, gracias. \loesung{¿Y tú?}
    \item B: \loesung{Me llamo [Name].}
    \item B: \loesung{¡Hasta luego! / ¡Adiós!}
\end{itemize}

\textbf{Aufgabe 3 (ser):} (4P)
\begin{itemize}
    \item yo \loesung{soy}, tú \loesung{eres}, él/ella \loesung{es}, nosotros \loesung{somos}
\end{itemize}

\textbf{Aufgabe 4 (Artikel):} (6P)
\begin{itemize}
    \item a) \loesung{el} libro, b) \loesung{la} mesa, c) \loesung{los} chicos
    \item d) \loesung{la} profesora, e) \loesung{las} sillas, f) \loesung{el} cuaderno
\end{itemize}

\section*{6. Stundenleistung SL-01d -- Übersicht}

\begin{tabular}{|c|l|c|c|}
\hline
\textbf{Aufg.} & \textbf{Inhalt} & \textbf{Punkte} & \textbf{Diff.} \\
\hline
1 & Zuordnung Begrüßungen & 4 & alle \\
\hline
2 & ser-Konjugation & 4 & alle \\
\hline
3 & Artikel einsetzen & 6 & alle \\
\hline
4 & Dialog vervollständigen & 6 & alle \\
\hline
\multicolumn{2}{|r|}{\textbf{Gesamt:}} & \textbf{20} & \\
\hline
\end{tabular}

\vspace{1em}

\textbf{Notenschlüssel (Klasse 7):}

\begin{tabular}{|c|c|c|c|c|c|c|}
\hline
\textbf{Note} & 1 & 2 & 3 & 4 & 5 & 6 \\
\hline
\textbf{ab \%} & 96 & 80 & 60 & 40 & 20 & <20 \\
\hline
\textbf{ab Punkte} & 19 & 16 & 12 & 8 & 4 & <4 \\
\hline
\end{tabular}

\section*{7. Reflexionsnotizen (nach Durchführung)}

\begin{tabularx}{\textwidth}{|l|X|}
\hline
\textbf{SL-Ergebnisse} & Ø: \rule{2cm}{0.4pt} / Beste: \rule{1cm}{0.4pt} / Schlechteste: \rule{1cm}{0.4pt} \\[1.5em]
\hline
\textbf{Problematische Aufgaben} & \\[2em]
\hline
\textbf{Für Sequenz 2} & \\[2em]
\hline
\end{tabularx}

\end{document}
